\chapter{Tuberculosis Skin Test (PPD or Mantoux Test)}

\section*{What Is This Test or Procedure?}
The tuberculosis skin test is an intradermal test that checks whether the immune system has reacted to proteins associated with \textit{Mycobacterium tuberculosis}.

\section*{Why It Is Performed}
It is used to screen people at risk for tuberculosis infection, including some people with exposure, healthcare work, residence in high-risk settings, or immune suppression. It does not by itself distinguish latent infection from active tuberculosis.

\section*{How It Is Performed}
A small amount of purified protein derivative (PPD) is injected just under the skin of the forearm, creating a small raised area. The injection site must be examined by a trained clinician 48 to 72 hours later; the diameter of firm swelling, not redness alone, is measured.

\section*{Preparation and Limitations}
No fasting or special preparation is needed. Tell the clinician about prior tuberculosis infection, BCG vaccination, immune suppression, or a previous severe reaction. A positive result may require further medical evaluation, and a negative result does not always exclude infection.

\section*{Risks and Results}
The test may cause temporary itching or soreness. Rarely, a stronger local reaction occurs. Interpretation depends on the size of induration and the person's risk factors; chest imaging and other evaluations may be needed when clinically indicated.

\section*{References}
\begin{enumerate}
  \item Centers for Disease Control and Prevention. Tuberculin skin testing. 2025.
  \item Current Medical Diagnosis and Treatment. McGraw-Hill; 2025.
\end{enumerate}
\clearpage
